% Part: first-order-logic
% Chapter: syntax-and-semantics
% Section: assignments

\documentclass[../../../include/open-logic-section]{subfiles}

\begin{document}

\olfileid{fol}{syn}{ass}

\olsection{Bedelingen}

\begin{explain}
Een !!{variable} bedeling~$s$ kent aan \emph{elke} variabele een
waarde toe, en daarvan zijn er oneindig veel. Strikt genomen is dat
niet nodig. We eisen alleen omwille van de eenvoud dat bedelingen aan
alle !!{variable}s een waarde toekennen. De waarde van een term~$t$ en
de vraag of !!a{formula}~$!A$ in !!a{structure} onder~$s$ wordt
vervuld, hangen uitsluitend af van wat de bedeling~$s$ toekent aan de
!!{variable}s in~$t$ en aan de vrije !!{variable}s van~$!A$. Dat is de
inhoud van de volgende twee proposities. We maken ‘hangt af van’
precies door te bewijzen dat twee bedelingen die op alle variabelen
in~$t$ overeenstemmen dezelfde waarde geven, en dat $!A$ onder de ene
bedeling wordt vervuld dan en slechts dan als zij onder de andere
wordt vervuld wanneer beide bedelingen op alle vrije variabelen
van~$!A$ overeenstemmen.
\end{explain}

\begin{prop}\ollabel{prop:valindep}
Als de !!{variable}s in een term~$t$ tot $x_1$, \dots,~$x_n$ behoren
en $s_1(x_i) = s_2(x_i)$ voor $i = 1$, \dots,~$n$, dan geldt
$\Value{t}{M}[s_1] = \Value{t}{M}[s_2]$.
\end{prop}

\begin{proof}
We bewijzen dit met inductie naar de complexiteit van~$t$. In het
basisgeval is~$t$ !!a{constant} of een van de variabelen~$x_1$,
\dots,~$x_n$. Als $t = c$, dan is
$\Value{t}{M}[s_1] = \Assign{c}{M} = \Value{t}{M}[s_2]$. Als
$t = x_i$, dan geldt volgens de aanname van de propositie
$s_1(x_i) = s_2(x_i)$, en dus
$\Value{t}{M}[s_1] = s_1(x_i) = s_2(x_i) = \Value{t}{M}[s_2]$.

Neem voor de inductiestap aan dat
$t = \Atom{f}{t_1, \dots, t_k}$ en dat de bewering geldt voor
$t_1$, \dots, $t_k$. Dan
\begin{align*}
  \Value{t}{M}[s_1] & = \Value{\Atom{f}{t_1, \dots, t_k}}{M}[s_1] \\
  & = \Assign{f}{M}(\Value{t_1}{M}[s_1], \dots, \Value{t_k}{M}[s_1]).
\intertext{Voor $j = 1$, \dots,~$k$ behoren de !!{variable}s van~$t_j$
    tot $x_1$, \dots,~$x_n$. Volgens de inductiehypothese is
    $\Value{t_j}{M}[s_1] = \Value{t_j}{M}[s_2]$. Bijgevolg}
\Value{t}{M}[s_1] & = \Value{\Atom{f}{t_1, \dots, t_k}}{M}[s_1] \\
  & = \Assign{f}{M}(\Value{t_1}{M}[s_1], \dots, \Value{t_k}{M}[s_1]) \\
  & = \Assign{f}{M}(\Value{t_1}{M}[s_2], \dots, \Value{t_k}{M}[s_2]) \\
  & = \Value{\Atom{f}{t_1, \dots, t_k}}{M}[s_2] = \Value{t}{M}[s_2].
\end{align*}
\end{proof}

\begin{prop}\ollabel{prop:satindep}
Als de vrije !!{variable}s in~$!A$ tot $x_1$, \dots,~$x_n$ behoren
en $s_1(x_i) = s_2(x_i)$ voor $i = 1$, \dots,~$n$, dan geldt
$\Sat{M}{!A}[s_1]$ dan en slechts dan als $\Sat{M}{!A}[s_2]$.
\end{prop}

\begin{proof}
We gebruiken inductie naar de complexiteit van~$!A$. In het
basisgeval is~$!A$ atomair en heeft~$!A$ een van de vormen:
\iftag{prvTrue}{$\ltrue$,}{}
\iftag{prvFalse}{$\lfalse$,}{}
$\Atom{R}{t_1, \dots, t_k}$ voor een $k$-plaatsig predicaat~$R$ en
termen $t_1$, \dots,~$t_k$, of $\eq[t_1][t_2]$ voor termen $t_1$
en~$t_2$. In de laatste twee gevallen bewijzen we alleen de voorwaartse
richting van de !!{biconditional}; de omgekeerde richting is
symmetrisch.

\begin{enumerate}
\tagitem{prvTrue}{%
  \indcase{!A}{\ltrue}{zowel $\Sat{M}{\indfrm}[s_1]$ als
    $\Sat{M}{\indfrm}[s_2]$ geldt.}}{}

\tagitem{prvFalse}{%
  \indcase{!A}{\lfalse}{zowel $\Sat/{M}{!A}[s_1]$ als
    $\Sat/{M}{!A}[s_2]$ geldt.}}{}

\item
  \indcase{!A}{\Atom{R}{t_1, \ldots, t_k}}{stel dat
    $\Sat{M}{\indfrm}[s_1]$. Dan
    \[
    \langle \Value{t_1}{M}[s_1], \ldots, \Value{t_k}{M}[s_1] \rangle
    \in \Assign{R}{M}.
    \]
    Voor $i = 1$, \dots,~$k$ is volgens \olref{prop:valindep}
    $\Value{t_i}{M}[s_1] = \Value{t_i}{M}[s_2]$. Daarom geldt ook
    $\langle \Value{t_1}{M}[s_2], \ldots, \Value{t_k}{M}[s_2] \rangle
    \in \Assign{R}{M}$, en dus $\Sat{M}{\indfrm}[s_2]$.}
\item
  \indcase{!A}{\eq[t_1][t_2]}{stel dat $\Sat{M}{\indfrm}[s_1]$.
    Dan is $\Value{t_1}{M}[s_1] = \Value{t_2}{M}[s_1]$. Bijgevolg
    \begin{align*}
      \Value{t_1}{M}[s_2] & = \Value{t_1}{M}[s_1]
      & \text{(volgens \olref{prop:valindep})} \\
      & = \Value{t_2}{M}[s_1] 
      & \text{(omdat $\Sat{M}{\eq[t_1][t_2]}[s_1]$)}\\
      &= \Value{t_2}{M}[s_2]
      & \text{(volgens \olref{prop:valindep}),}
    \end{align*}
    zodat $\Sat{M}{\eq[t_1][t_2]}[s_2]$.}
\end{enumerate}

Neem nu aan dat $\Sat{M}{!B}[s_1]$ dan en slechts dan als
$\Sat{M}{!B}[s_2]$ geldt voor elke !!{formula}~$!B$ die minder complex
is dan~$!A$. In de inductiestap onderscheiden we gevallen naar de
hoofdoperator van~$!A$. In elk geval bewijzen we alleen de
voorwaartse richting van de !!{biconditional}; de omgekeerde richting
is symmetrisch. Behalve bij de kwantoren passen we de
inductiehypothese toe op sub-!!{formula}s~$!B$ van~$!A$. De vrije
variabelen van~$!B$ behoren tot die van~$!A$. Als $s_1$ en $s_2$ op de
vrije variabelen van~$!A$ overeenstemmen, stemmen ze dus ook op die
van~$!B$ overeen en is de inductiehypothese op~$!B$ van toepassing.

\begin{enumerate}
\tagitem{defNot}{}{%
  \iftag{probNot}{%
    \indcase!{!A}{\lnot !B}{}}{%
    \indcase{!A}{\lnot !B}{als $\Sat{M}{\indfrm}[s_1]$, dan
      $\Sat/{M}{!B}[s_1]$. Volgens de inductiehypothese geldt dus
      $\Sat/{M}{!B}[s_2]$, en bijgevolg $\Sat{M}{\indfrm}[s_2]$.}}}

\tagitem{defAnd}{}{%
  \iftag{probAnd}{%
    \indcase!{!A}{!B \land !C}{}}{%
    \indcase{!A}{!B \land !C}{als $\Sat{M}{\indfrm}[s_1]$, dan
      $\Sat{M}{!B}[s_1]$ en $\Sat{M}{!C}[s_1]$. Volgens de
      inductiehypothese gelden $\Sat{M}{!B}[s_2]$ en
      $\Sat{M}{!C}[s_2]$. Bijgevolg
      $\Sat{M}{\indfrm}[s_2]$.}}}

\tagitem{defOr}{}{%
  \iftag{probOr}{%
    \indcase!{!A}{!B \lor !C}{}}{%
    \indcase{!A}{!B \lor !C}{als $\Sat{M}{\indfrm}[s_1]$, dan
      $\Sat{M}{!B}[s_1]$ of $\Sat{M}{!C}[s_1]$. Volgens de
      inductiehypothese geldt $\Sat{M}{!B}[s_2]$ of
      $\Sat{M}{!C}[s_2]$, en dus $\Sat{M}{\indfrm}[s_2]$.}}}

\tagitem{defIf}{}{%
  \iftag{probIf}{%
    \indcase!{!A}{!B \lif !C}{}}{%
    \indcase{!A}{!B \lif !C}{als $\Sat{M}{\indfrm}[s_1]$, dan
    $\Sat/{M}{!B}[s_1]$ of $\Sat{M}{!C}[s_1]$. Volgens de
    inductiehypothese geldt $\Sat/{M}{!B}[s_2]$ of
    $\Sat{M}{!C}[s_2]$, en dus $\Sat{M}{\indfrm}[s_2]$.}}}

\tagitem{defIff}{}{%
  \iftag{probIff}{%
    \indcase!{!A}{!B \liff !C}{}}{%
    \indcase{!A}{!B \liff !C}{als $\Sat{M}{\indfrm}[s_1]$, dan gelden
      óf $\Sat{M}{!B}[s_1]$ en $\Sat{M}{!C}[s_1]$, óf
      $\Sat/{M}{!B}[s_1]$ en $\Sat/{M}{!C}[s_1]$. Volgens de
      inductiehypothese gelden óf $\Sat{M}{!B}[s_2]$ en
      $\Sat{M}{!C}[s_2]$, óf $\Sat/{M}{!B}[s_2]$ en
      $\Sat/{M}{!C}[s_2]$. In beide gevallen geldt
      $\Sat{M}{\indfrm}[s_2]$.}}}

\tagitem{defEx}{}{%
  \iftag{probEx}{%
    \indcase!{!A}{\lexists[x][!B]}{}}{%
    \indcase{!A}{\lexists[x][!B]}{als $\Sat{M}{\indfrm}[s_1]$, dan is
      er een $m \in \Domain{M}$ waarvoor
      $\Sat{M}{!B}[\Subst{s_1}{m}{x}]$. Stel
      $s_1' = \Subst{s_1}{m}{x}$ en $s_2' =
      \Subst{s_2}{m}{x}$. De vrije variabelen van~$!B$ behoren tot
      $x_1$, \dots, $x_n$ en~$x$. Er geldt
      $s_1'(x_i) = s_2'(x_i)$: $s_1'$ en $s_2'$ zijn respectievelijk
      $x$-varianten van $s_1$ en~$s_2$, en volgens de aanname is
      $s_1(x_i) = s_2(x_i)$. Door de definitie is bovendien
      $s_1'(x) = s_2'(x) = m$. De aldus gedefinieerde bedelingen
      $s_1'$ en~$s_2'$ voldoen dus aan de voorwaarden. De
      inductiehypothese is van toepassing op $!B$, $s_1'$ en~$s_2'$, zodat
      $\Sat{M}{!B}[s_2']$. Omdat $s_2' = \Subst{s_2}{m}{x}$, is er
      derhalve een $m \in \Domain{M}$ waarvoor
      $\Sat{M}{!B}[\Subst{s_2}{m}{x}]$, en dus
      $\Sat{M}{\indfrm}[s_2]$.}}}

\tagitem{defAll}{}{%
  \iftag{probAll}{%
    \indcase!{!A}{\lforall[x][!B]}{}}{%
    \indcase{!A}{\lforall[x][!B]}{als $\Sat{M}{\indfrm}[s_1]$, dan
      geldt voor elke $m \in \Domain{M}$:
      $\Sat{M}{!B}[\Subst{s_1}{m}{x}]$. We moeten aantonen dat ook voor
      elke $m \in \Domain{M}$ geldt:
      $\Sat{M}{!B}[\Subst{s_2}{m}{x}]$. Kies dus een willekeurig
      $m \in \Domain{M}$ en stel
      $s_1' = \Subst{s_1}{m}{x}$ en $s_2' =
      \Subst{s_2}{m}{x}$. Dan geldt $\Sat{M}{!B}[s_1']$. De vrije
      variabelen van~$!B$ behoren tot $x_1$, \dots, $x_n$ en~$x$.
      Er geldt $s_1'(x_i) = s_2'(x_i)$, omdat $s_1'$ en $s_2'$
      respectievelijk $x$-varianten van $s_1$ en~$s_2$ zijn en volgens
      de aanname $s_1(x_i) = s_2(x_i)$. Door hun definitie is ook
      $s_1'(x) = s_2'(x) = m$. De aldus gedefinieerde bedelingen
      $s_1'$ en~$s_2'$ voldoen dus aan de voorwaarden. De
      inductiehypothese is van toepassing op~$!B$, $s_1'$ en~$s_2'$, en levert
      $\Sat{M}{!B}[s_2']$. Dit geldt voor elk $m \in \Domain{M}$, dat
      wil zeggen: $\Sat{M}{!B}[\Subst{s_2}{m}{x}]$ voor alle
      $m \in \Domain{M}$. Dus $\Sat{M}{\indfrm}[s_2]$.}}}
\end{enumerate}
Uit de inductie volgt dat $\Sat{M}{!A}[s_1]$ dan en slechts dan als
$\Sat{M}{!A}[s_2]$, zodra de vrije !!{variable}s in~$!A$ tot
$x_1$, \dots, $x_n$ behoren en $s_1(x_i)=s_2(x_i)$ voor
$i = 1$, \dots,~$n$.
\end{proof}

\begin{probtag}{probNot,probOr,probAnd,probIf,probIff,probEx,probAll}
Voltooi het bewijs van \olref[fol][syn][ass]{prop:satindep}.
\end{probtag}

\begin{explain}
!!^{sentence}s hebben geen vrije variabelen. Twee willekeurige
bedelingen kennen dus aan alle nul vrije variabelen van een zin
dezelfde objecten toe. Uit de zojuist bewezen propositie volgt dat de
vraag of !!a{sentence} in een structuur onder een bedeling wordt
vervuld, volledig onafhankelijk is van die bedeling. We leggen dit
feit vast. Het rechtvaardigt de daaropvolgende definitie van vervulling
van !!a{sentence} in !!a{structure}, waarin geen bedeling wordt genoemd.
\end{explain}

\begin{cor}
\ollabel{cor:sat-sentence}
Als $!A$ !!a{sentence} is en $s$ een bedeling, dan geldt
$\Sat{M}{!A}[s]$ dan en slechts dan als $\Sat{M}{!A}[s']$ voor elke
bedeling~$s'$.
\end{cor}

\begin{proof}
  Kies een willekeurige bedeling~$s'$. Omdat $!A$ !!a{sentence} is,
  heeft zij geen vrije variabelen. Elke bedeling~$s'$ kent dus triviaal
  aan alle vrije variabelen van~$!A$ hetzelfde toe als~$s$. Aan de
  voorwaarde van \olref{prop:satindep} is bijgevolg voldaan, en
  $\Sat{M}{!A}[s]$ geldt dan en slechts dan als
  $\Sat{M}{!A}[s']$.
\end{proof}

\begin{defn}
\ollabel{defn:satisfaction}
Als $!A$ !!a{sentence} is, zeggen we dat !!a{structure}~$\Struct M$
$!A$ \emph{vervult}, genoteerd $\Sat{M}{!A}$, dan en slechts dan als
$\Sat{M}{!A}[s]$ voor alle bedelingen~$s$.
\end{defn}

Als $\Sat{M}{!A}$, zeggen we ook eenvoudigweg dat \emph{$!A$ waar is
in~$\Struct{M}$.} Het begrip vervulling breidt zich op natuurlijke
wijze uit van afzonderlijke !!{sentence}s tot verzamelingen
!!{sentence}s.

\begin{defn}
\ollabel{defn:sat}
Beschouw een verzameling~$\Gamma$ van !!{sentence}s; we blijven deze
verzameling aanduiden met~$\Gamma$. We zeggen dat
!!a{structure}~$\Struct M$ de verzameling~$\Gamma$ \emph{vervult},
genoteerd $\Sat{M}{\Gamma}$, dan en slechts dan als
$\Sat{M}{!A}$ voor alle $!A \in \Gamma$.
\end{defn}

\begin{prop}\ollabel{prop:sentence-sat-true}
  Zij $\Struct{M}$ !!a{structure}, zij $!A$ !!a{sentence} en zij $s$
  een bedeling. Dan geldt $\Sat{M}{!A}$ dan en slechts dan als
  $\Sat{M}{!A}[s]$.
\end{prop}

\begin{proof}
Opgave.
\end{proof}

\begin{prob}
Bewijs \olref[fol][syn][ass]{prop:sentence-sat-true}.
\end{prob}

\begin{prop}\ollabel{prop:sat-quant}
  Stel dat $!A(x)$ alleen $x$ als vrije variabele bevat en dat
  $\Struct M$ !!a{structure} is. Dan geldt:
  \begin{tagenumerate}{prvEx,prvAll}
    \tagitem{prvEx}{$\Sat{M}{\lexists[x][!A(x)]}$ dan en slechts dan als
      $\Sat{M}{!A(x)}[s]$
      voor ten minste één bedeling~$s$.}{}
    \tagitem{prvAll}{$\Sat{M}{\lforall[x][!A(x)]}$ dan en slechts dan als
      $\Sat{M}{!A(x)}[s]$
  voor alle bedelingen~$s$.}{}
\end{tagenumerate}
\end{prop}

\begin{proof}
Opgave.
\end{proof}

\begin{prob}
Bewijs \olref[fol][syn][ass]{prop:sat-quant}.
\end{prob}

\begin{prob}
\DeclareRobustCommand{\VDash}{\mathrel{||}\joinrel\Relbar}  
Stel dat $\Lang L$ een taal zonder !!{function}s is. Gegeven
!!a{structure}~$\Struct M$, !!a{constant}~$c$ en
$a \in \Domain M$, definiëren we $\Struct M[a/c]$ als de
!!{structure} die gelijk is aan~$\Struct M$, behalve dat
$\Assign{c}{M[a/c]} = a$. Definieer de relatie
$\Struct M \VDash !A$ voor !!{sentence}s~$!A$ als volgt:
\begin{enumerate}
\tagitem{prvFalse}{%
  \indcase{!A}{\lfalse}{niet $\Struct{M} \VDash \indfrm$.}}{}

\tagitem{prvTrue}{%
  \indcase{!A}{\ltrue}{$\Struct{M} \VDash \indfrm$.}}{}

\item \indcase{!A}{\Atom{R}{d_1, \dots, d_n}}{$\Struct M \VDash \indfrm$
  dan en slechts dan als
  $\langle \Assign{d_1}{M}, \dots, \Assign{d_n}{M} \rangle \in
  \Assign{R}{M}$.}

\item \indcase{!A}{\eq[d_1][d_2]}{$\Struct M \VDash \indfrm$ dan en
  slechts dan als $\Assign{d_1}{M} = \Assign{d_2}{M}$.}

\tagitem{prvNot}{%
  \indcase{!A}{\lnot !B}{$\Struct{M} \VDash \indfrm$ dan en slechts
    dan als niet $\Struct{M} \VDash {!B}$.}}{}

\tagitem{prvAnd}{%
  \indcase{!A}{(!B \land !C)}{$\Struct{M} \VDash
    \indfrm$ dan en slechts dan als $\Struct{M} \VDash!B$ en
    $\Struct{M} \VDash !C$.}}{}

\tagitem{prvOr}{%
  \indcase{!A}{(!B \lor !C)}{$\Struct{M} \VDash \indfrm$ dan en
    slechts dan als $\Struct{M} \VDash !B$ of
    $\Struct{M} \VDash !C$ (of beide).}}{}

\tagitem{prvIf}{%
  \indcase{!A}{(!B \lif !C)}{$\Struct{M} \VDash \indfrm$ dan en
    slechts dan als niet $\Struct {M} \VDash !B$ of
    $\Struct M \VDash !C$ (of beide).}}{}

\tagitem{prvIff}{%
  \indcase{!A}{(!B \liff !C)}{$\Struct{M} \VDash {\indfrm}$ dan en
    slechts dan als óf zowel $\Struct{M} \VDash {!B}$ als
    $\Struct{M} \VDash {!C}$ geldt, óf noch
    $\Struct{M} \VDash {!B}$ noch $\Struct{M} \VDash {!C}$ geldt.}}{}

\tagitem{prvAll}{%
  \indcase{!A}{\lforall[x][!B]}{$\Struct{M} \VDash {\indfrm}$ dan en
    slechts dan als voor alle $a \in \Domain{M}$ geldt
    $\Struct{M[a/c]} \VDash \Subst{!B}{c}{x}$, mits $c$ niet
    in~$!B$ voorkomt.}}{}

\tagitem{prvEx}{%
  \indcase{!A}{\lexists[x][!B]}{$\Struct{M} \VDash
  {\indfrm}$ dan en slechts dan als er een $a \in \Domain M$ is
  waarvoor $\Struct{M[a/c]} \VDash \Subst{!B}{c}{x}$, mits $c$ niet
  in~$!B$ voorkomt.}}{}
\end{enumerate}
Laat $x_1$, \dots, $x_n$ alle vrije !!{variable}s in~$!A$ zijn,
laat $c_1$, \dots, $c_n$ constantsymbolen zijn die niet in~$!A$
voorkomen, neem $a_1$, \dots, $a_n \in \Domain M$ en stel
$s(x_i) = a_i$.

Toon aan dat $\Sat{M}{!A}[s]$ dan en slechts dan als
$\Struct M[a_1/c_1,\dots,a_n/c_n]
\VDash \Subst{\Subst{!A}{c_1}{x_1}\dots}{c_n}{x_n}$.

(Deze opgave laat zien dat men semantiek voor de eerste-ordelogica kan
geven zonder bedelingen te gebruiken.)
\end{prob}

\begin{prob}
Stel dat $f$ een functiesymbool is dat niet in~$!A(x,y)$ voorkomt.
Toon aan dat er !!a{structure}~$\Struct{M}$ bestaat waarvoor
$\Sat{M}{\lforall[x][\lexists[y][!A(x,y)]]}$ dan en slechts dan als er
!!a{structure}~$\Struct M'$ bestaat waarvoor
$\Sat{M'}{\lforall[x][!A(x,f(x))]}$.

(Deze opgave is een bijzonder geval van de stelling van Skolem;
$\lforall[x][!A(x,f(x))]$ heet een \emph{Skolemnormaalvorm} van
$\lforall[x][\lexists[y][!A(x,y)]]$.)
\end{prob}

\end{document}
